\clearpage
\phantomsection
\addcontentsline{toc}{chapter}{原课程信息与课历}
\markboth{原课程信息与课历}{原课程信息与课历}
\begin{center}
{\color{structurecolor}\bfseries\fontsize{14.4}{20}\selectfont 原课程信息与课历\par}
\end{center}
\noindent\textbf{课程身份.} MIT OpenCourseWare 18.950, \textit{Differential Geometry}, Fall 2008; 原授课教师 Paul Seidel, 数学系, 本科课程. 以下课程信息取自\href{https://ocw.mit.edu/courses/18-950-differential-geometry-fall-2008/pages/syllabus/}{官方教学大纲}, 不把本整理本的教学安排冒称原课历.

\noindent\textbf{课程目标与先修.} 原课程以曲率为中心, 采用定义、定理和证明的严格讲授, 同时强调具体几何. 官方先修为 Analysis I (18.100), 加 Linear Algebra (18.06 或18.700) 或 Algebra I (18.701). 学习者应熟悉多元微积分和线性代数. 本册还在实际证明中调用光滑常微分方程、逆函数与隐函数定理及 Green 公式; 这些是整理本的具体先修说明, 不是另造官方选课门槛.

\noindent\textbf{授课时间与讲次.} 公开记录只确定授课学期为2008年秋季; 现存课程页面未公开每周授课星期和时刻, 也未提供逐日 Calendar. 因而下表保留\href{https://ocw.mit.edu/courses/18-950-differential-geometry-fall-2008/pages/lecture-notes/}{官方讲义目录}的讲次顺序, 不添加推测日期.
\begin{center}
\begin{tabular}{p{.15\textwidth}p{.72\textwidth}}
\toprule
讲次 & 官方四章的内容\\
\midrule
1--10 & 平面曲线的局部与整体几何.\\
11--23 & 超曲面的局部几何.\\
24--35 & 超曲面的整体几何.\\
36--41 & 长度与距离的几何.\\
\bottomrule
\end{tabular}
\end{center}
\noindent\textbf{考核.} 官方大纲规定作业占30\%, 期中考试占30\%, 期末考试占40\%. 这说明原课程设有两类考试, 不说明其试卷已公开. 现存公开目录提供10份 Homework PDF, 没有公开的官方解答或考试原卷; 本册不编造考试题面. 作业截止时间仅按各份原件实际印出的信息记载.

\noindent\textbf{指定教材与材料边界.} 官方指定 Wolfgang K\"uhnel, \textit{Differential Geometry: Curves--Surfaces--Manifolds}, AMS, 2002, 并说明大致跟随前半部. 官方还列有 Spivak、do Carmo、Pressley 等参考书. 本项目没有取得并复制这些商业教材正文. 作业6.1提到教材中的翻转隆起方法, 本册另作自足构造, 不冒称重现未读教材. 作业3.3和9.4所引两处课堂编号未能在公开修订讲义中定位, 保留原指涉并标明缺口.

\noindent\textbf{正文与作业的关系.} 先保留已交付的12章基础正文和36道练习解答, 再按原次序并入10份官方作业, 共32个顶层题目. 四道题另有10个带标号子问, 不将它们重复加到32题中. 可从公开材料恢复的要求均给中文题面和独立解答; 作业3.3最后一个指涉及9.4整题仍有上述来源缺口. 能与基础册准确匹配的原题注明映射并复用相应证明; 仅主题相近者不冒称同一道题. 新增解答均为AI独立编写与复审, 不是MIT或原教师发布的官方答案. 10份作业的主要整理位置为本微分几何册; 拓扑册只需互链, 不重复制作另一套答案.

\noindent\textbf{许可与署名.} 原课程作者、MIT机构、Fall 2008年份和官方来源保留. 依MIT OCW当前默认CC BY-NC-SA 4.0许可非商业、相同方式分享; 单独权利声明优先. 独立解答、中文翻译和章节重组属于本整理本的改动, 不表示MIT或Paul Seidel认可本项目.
\clearpage
